
\documentclass{article}
\usepackage{colortbl}
\usepackage{makecell}
\usepackage{multirow}
\usepackage{supertabular}

\begin{document}

\newcounter{utterance}

\centering \large Interaction Transcript for game `cladder', experiment `full\_v1.5\_default', episode 4257 with Qwen3.6{-}27B.
\vspace{24pt}

{ \footnotesize  \setcounter{utterance}{1}
\setlength{\tabcolsep}{0pt}
\begin{supertabular}{c@{$\;$}|p{.15\linewidth}@{}p{.15\linewidth}p{.15\linewidth}p{.15\linewidth}p{.15\linewidth}p{.15\linewidth}}
    \# & \multicolumn{2}{c}{Player} && \multicolumn{2}{c}{Game Master} \\
    \hline

    \theutterance \stepcounter{utterance}  
    & & & \multicolumn{4}{p{0.6\linewidth}}{
        \cellcolor[rgb]{0.9,0.9,0.9}{
            \makecell[{{p{\linewidth}}}]{
                \texttt{\tiny{[P1$\langle$GM]}}
                \texttt{You are an expert in causal inference. The following question is not a typical commonsense query, but rather a meticulously designed question created by a professor specializing in causal inference, intended to assess the students' mastery of the course content.} \\
\texttt{Imagine a self{-}contained, hypothetical world with only the following conditions, and without any unmentioned factors or causal relationships: Kidney stone size has a direct effect on treatment and lip thickness. Treatment has a direct effect on lip thickness. The overall probability of receives treatment is 48\%. For patients not receiving treatment, the probability of thick lips is 64\%. For patients receiving treatment, the probability of thick lips is 29\%. Is thick lips more likely than thin lips overall?} \\
\texttt{Start your answer with 'Yes' or 'No', followed by additional reasoning or evidence to support your explanation.} \\
            }
        }
    }
     \\ \\

    \theutterance \stepcounter{utterance}  
    & \multicolumn{4}{p{0.6\linewidth}}{
        \cellcolor[rgb]{0.561,0.737,0.561}{
            \makecell[{{p{\linewidth}}}]{
                \texttt{\tiny{[P1$\rangle$GM]}}
                \texttt{No} \\
\\ 
\texttt{To determine whether thick lips are more likely than thin lips overall, we must calculate the total probability of having thick lips, denoted as \$P(L=1)\$, where \$L=1\$ represents thick lips and \$L=0\$ represents thin lips.} \\
\\ 
\texttt{According to the law of total probability, we can decompose \$P(L=1)\$ based on the treatment variable \$T\$:} \\
\texttt{\$\$P(L=1) = P(L=1 | T=1)P(T=1) + P(L=1 | T=0)P(T=0)\$\$} \\
\\ 
\texttt{From the problem description, we have the following values:} \\
\texttt{1.  The overall probability of receiving treatment: \$P(T=1) = 0.48\$.} \\
\texttt{2.  Consequently, the probability of not receiving treatment: \$P(T=0) = 1 {-} P(T=1) = 1 {-} 0.48 = 0.52\$.} \\
\texttt{3.  The probability of thick lips given no treatment: \$P(L=1 | T=0) = 0.64\$.} \\
\texttt{4.  The probability of thick lips given treatment: \$P(L=1 | T=1) = 0.29\$.} \\
\\ 
\texttt{Now, we substitute these values into the equation:} \\
\texttt{\$\$P(L=1) = (0.29 \textbackslash{}times 0.48) + (0.64 \textbackslash{}times} \\
            }
        }
    }
    & & \\ \\

    \theutterance \stepcounter{utterance}  
    & & & \multicolumn{2}{p{0.3\linewidth}}{
        \cellcolor[rgb]{0.9,0.9,0.9}{
            \makecell[{{p{\linewidth}}}]{
                \texttt{\tiny{[GM$|$GM]}}
                \texttt{No} \\
            }
        }
    }
    & & \\ \\

    \theutterance \stepcounter{utterance}  
    & & & \multicolumn{2}{p{0.3\linewidth}}{
        \cellcolor[rgb]{0.9,0.9,0.9}{
            \makecell[{{p{\linewidth}}}]{
                \texttt{\tiny{[GM$|$GM]}}
                \texttt{no} \\
            }
        }
    }
    & & \\ \\

    \theutterance \stepcounter{utterance}  
    & & & \multicolumn{2}{p{0.3\linewidth}}{
        \cellcolor[rgb]{0.9,0.9,0.9}{
            \makecell[{{p{\linewidth}}}]{
                \texttt{\tiny{[GM$|$GM]}}
                \texttt{game\_result = WIN} \\
            }
        }
    }
    & & \\ \\

\end{supertabular}
}

\end{document}
